\section{Введение}

Основной операцией чистого логического программирования является синтаксическая унификация~\cite{robinson1965machine} термов первого порядка (они же эрбрановы деревья). Унификация позволяет решать уравнения между термами за счёт уточнения их свободных переменных. Синтаксическая унификация отличается тем, что любые функциональные символы полагаются неинтерпретируемыми -- то есть, применение функционального символа к термам $f\,(t_1, ..., t_n)$ интерпретируется как идентичная синтаксическая конструкция. Все такие конечные синтаксические записи образуют эрбранову вселенную, адекватность применения которой обосновывается теоремой Эрбрана~\cite{herbrand1930}.

Помимо логического программирования, синтаксическая унификация находит применение при выводе и проверке типов в языках программирования. В зависимости от конкретной системы типов, иногда полезно расширить пространство возможных синтаксических форм бесконечными -- это позволяет говорить о рекурсивных типах (например, в языке программирования \OCaml~\cite{gauthier2004numbering}). Тем не менее, произвольно-бесконечные термы не имеют практической применимости в вычислительной математической логике, потому что для их порождения потребовалось бы затратить бесконечное время работы программы. Вместо этого рассматривают более узкий класс возможно-бесконечных термов -- рациональные термы. Рациональным термом~\cite{colmerauer1982prolog} называют такой возможно-бесконечный терм, множество подтермов (в смысле поддеревьев) которого конечно. В частности, любой рекурсивный тип можно представить как рациональный терм.

В логическом программировании унификация рациональных термов также находит применение, но не столько за её выразительную способность, сколько за скорость работы. Проблема в том, что для ограничения конечности термов требуется выполнение трудоёмкой\footnote{Например, реализация конкатенации двух списков при использовании классического алгоритма унификации требует затрат $\mathcal{O}(n^2)$ времени вместо $\mathcal{O}(n)$~\cite{colmerauer1982prolog}.} <<проверки вхождения>> (англ. <<occurs check>>) или применение более сложных алгоритмов унификации~\cite{martelli1982efficient}, обладающих непрактичными накладными расходами времени работы. Из-за этого практические реализации языков логического программирования (например, SWI-Prolog) используют унификацию рациональных термов, позволяя пользователю в частных случаях проверять отсутствие зацикливания внутри термов.

Несмотря на то, что рациональная унификация является хорошо изученной темой, в рамках которой было предложено множество алгоритмов~\cite{huet1976resolution,mukai1983unification,martelli1984efficient,jaffar1984efficient}, обширное применение получил только первый из перечисленных алгоритм Юэ~\cite{huet1976resolution}. Оптимальная реализация данного алгоритма действительно хорошо работает в существующих приложениях за счёт использования изменяемой памяти. В реализациях языка \Prolog это возможно благодаря использованию WAM (Warren Abstract Machine~\cite{warren1983abstract}), которая реализует поиск в глубину с откатами -- то есть, при ошибке в активной ветке дерева поиска выполняется откат состояний изменяемых переменных в памяти до ближайшего ветвления (backtracking). А в выводе типов основной упор делается на детерминированность поиска -- то есть, если произошла ошибка, нет смысла рассматривать никакие альтернативные ветки.

Тем не менее, в логическом программировании выделяют отдельную область реляционного программирования и язык \miniKanren~\cite{friedman2005reasoned}. Одним из отличий \miniKanren от \Prolog является использование поиска с перестановками (interleaving search~\cite{kiselyov2005backtracking}) вместо поиска в глубину. Такой выбор обусловлен полнотой поиска -- способностью находить все возможные решения за некоторое конечное время, -- в отличие от поиска в глубину, который может уйти в непродуктивный бесконечный цикл. Для реализации поиска с перестановками обычно используют персистентные структуры данных для представления отдельных состояний в дереве поиска, потому что процедура поиска вынуждена <<одновременно>> поддерживать несколько <<активных>> состояний, переключение между которыми может выполняться неожиданным образом.

В этой работе мы описываем собственный алгоритм для унификации рациональных термов на основе подхода Мартелли и Росси~\cite{martelli1984efficient}. Наш алгоритм обладает доказанной корректностью и завершаемостью, формально верифицированными в системе автоматической проверки доказательств \Rocq. Сами доказательства не являются предметом данной статьи, но в тексте мы иногда ссылаемся на конкретные определения из реализации на \Rocq\footnote{\url{https://github.com/dboulytchev/miniKanren-coq/tree/master/src/Rational}}.

Основным результатом этой работы является экспериментальный анализ производительности нашего алгоритма и алгоритма Юэ~\cite{huet1976resolution} в сравнении с классическим алгоритмом Робинсона с треугольными подстановками~\cite{baader2001unification}. Представленный анализ показывает, что в практических сценариях наш алгоритм обходит по скорости алгоритм Юэ будучи реализованным персистентно, хотя и уступает в эфемерной (не персистетной) реализации. В то же время нам удалось показать, что существуют сценарии, когда алгоритм Юэ работает значительно (более, чем константно) медленнее нашего.


\subsection{Рациональные термы}

Пусть заданы счётно-бесконечные множества переменных $x, y, z, ... \in \V$ и конструкторов (функциональных символов) $f, g, h, ... \in \C$, а также функция $\arity : \C \rightarrow \N$, определяющая арности конструкторов. Для удобства, мы используем запись $f^n \in \C$, где $f \in \C$, а $n = \arity(f)$.

Определим отображение $F : \text{Set} \rightarrow \text{Set}$, и множества термов $T = \mu F$ и (возможно-)бесконечных термов $T^\infty = \nu F$ как наименьшую и наибольшую~\cite{simon2006coinductive} неподвижные точки этого отображения соответственно:
\[F(X) = \V \cup \{ f^n\,(t_1, ..., t_n) \mid f^n \in \C, t_i \in X \}.\]
Множество подтермов $\subterms : T^\infty \rightarrow 2^{T^\infty}$ определяется индуктивно:
\begin{itemize}
    \item $t \in \subterms(t)$;
    \item $t \in \subterms(t_i) \implies t \in \subterms(f^n\,(..., t_i, ...))$.
\end{itemize}
Тогда множество рациональных термов определяется как:
\[T^R = \{ t \in T^\infty \mid \subterms(t) \text{ конечно} \}.\]
Естественно, что $T \subset T^R \subset T^\infty$, поэтому многие определения даются для $T^\infty$ и автоматически действуют для $T$ и $T^R$.

Подстановкой $\sigma \in \Sigma$ мы называем такое отображение $\V \rightarrow T^\infty$, которое обладает конечным доменом подстановки $\dom(\sigma) = \{ x \in \V \mid \sigma(x) \neq x \}$. Применением подстановки $\sigma$ к терму $t \in T^\infty$ называется гомоморфное расширение применения подстановки как функции и записывается в постфиксной форме как $t \sigma$:
\begin{equation*}
    \begin{array}{r@{{}={}}l}
        x \sigma & \sigma(x) \\
        f^n\,(..., t_i, ...) \sigma & f^n\,(..., t_i \sigma, ...) \\
    \end{array}
\end{equation*}
